\FloatBarrier
\section{KITTI: Frozen Detection and Detector Adaptation}
\label{sec:kitti-results}

\subsection{Three-Epoch Full-Validation Comparison}
\label{sec:kitti-final}

The initial three-epoch PU-GCN study completed all 20 planned detector evaluations,
each covering the same 3,769 validation frames. The matrix is intentionally
asymmetric: PointRCNN includes direct-input adaptation in both lines,
whereas CenterPoint adaptation was completed for observed-first inputs and
the sparse baseline. An unexecuted direct-adapted CenterPoint arm is not
inferred. Tables~\ref{tab:pr-a}--\ref{tab:cp-b} preserve this original matrix.
The later observed-first schedule extensions are reported separately in
Section~\ref{sec:adaptation-evidence}.

Throughout this section, \emph{frozen} denotes the released detector weights,
and \emph{adapted} denotes detector weights after training on the corresponding
representation. PU-GCN itself remained fixed at its released PU1K
\texttt{model-100} checkpoint. Direct inputs contain the strict network
output; observed-first inputs concatenate all observed rows with a
deterministically selected set of predicted rows. Those input constructions
have the same total cardinality, but different measured-point preservation.

\subsection{PointRCNN}
\label{sec:results-pointrcnn}

For PointRCNN, frozen direct upsampling reduced Line~A Car Moderate
3D AP$_{R40}$ from 81.95 to 61.93. Preserving observations increased it
to 64.27. Adaptation further raised the direct and observed-first results
to 68.88 and 71.01, respectively (Table~\ref{tab:pr-a}). Observation
preservation helped under both weight regimes, but the best three-epoch arm
remained 10.95 AP points below the original frozen baseline.

\begin{table}[t]\centering\setlength{\tabcolsep}{5pt}
\begin{tabular}{llrrrr}\toprule
Input & Weights & Easy & Moderate & Hard & $\Delta$ Mod. \\
\midrule\tableinput{tables/kitti_pointrcnn_A.tex}\bottomrule
\end{tabular}
\caption{PointRCNN Line~A Car 3D AP$_{R40}$, 3,769 frames per row.
All deltas reference the original frozen baseline. No adapted original
baseline was trained, so adapted-row deltas change both input and weights.}
\label{tab:pr-a}\end{table}

Line~B revealed a larger loss under direct insertion. The frozen sparse
baseline reached 65.58, compared with 30.09 for direct PU-GCN and 35.53
for observed-first PU-GCN. Training the detector on the direct representation
raised AP to 46.61, and observed-first adaptation raised it to 55.21
(Table~\ref{tab:pr-b}). This was substantial recovery relative to the
frozen generated inputs, but it did not exceed the adapted sparse baseline
at 68.33. The residual observed-first deficit was 13.12 AP points.

\begin{table}[t]\centering\setlength{\tabcolsep}{5pt}
\begin{tabular}{llrrrr}\toprule
Input & Weights & Easy & Moderate & Hard & $\Delta$ Mod. \\
\midrule\tableinput{tables/kitti_pointrcnn_B.tex}\bottomrule
\end{tabular}
\caption{PointRCNN Line~B Car 3D AP$_{R40}$. Every row evaluates 3,769
frames. Frozen rows reference the frozen sparse baseline; adapted rows
reference the separately adapted sparse baseline.}
\label{tab:pr-b}\end{table}

The sparse-baseline adaptation is an essential control. It improved
PointRCNN from 65.58 to 68.33 without generating points. Comparing adapted
upsampling only with a frozen sparse baseline would attribute some benefit
of downstream optimisation to the point generator. The paired rows in
Table~\ref{tab:pr-b} keep that control explicit. These results favour a representation mismatch
explanation for part of the frozen-model loss, while the remaining deficit
shows that the three-epoch adaptation procedure did not restore baseline
performance; it does not uniquely identify the remaining cause.

\subsection{CenterPoint and Class Dependence}
\label{sec:results-centerpoint}

CenterPoint recovered more of the original performance under observed-first
adaptation. Its Line~A baseline was 79.28 Car Moderate 3D AP$_{R40}$.
Frozen direct and observed-first inputs reached 60.61 and 63.21, while
adapted observed-first reached 74.70 (Table~\ref{tab:cp-a}). The latter
remained 4.58 AP points below the original frozen detector. As for
PointRCNN, there was no independently adapted original Line~A baseline;
this residual is a comparison between two complete systems rather than a
controlled effect of geometry at fixed weights.

\begin{table}[t]\centering\setlength{\tabcolsep}{5pt}
\begin{tabular}{llrrrr}\toprule
Input & Weights & Easy & Moderate & Hard & $\Delta$ Mod. \\
\midrule\tableinput{tables/kitti_centerpoint_A.tex}\bottomrule
\end{tabular}
\caption{CenterPoint Line~A Car 3D AP$_{R40}$, 3,769 frames per row.
Deltas reference the frozen original baseline. The direct-adapted
configuration was not part of the completed matrix.}\label{tab:cp-a}
\end{table}

In Line~B, frozen direct input reached 35.35 against the sparse
baseline's 64.60. Observation preservation increased it to 44.09, and
adaptation raised it to 61.66. The adapted sparse baseline reached 68.05,
leaving a matched-weight-family deficit of 6.39 AP points
(Table~\ref{tab:cp-b}). Greater recovery than PointRCNN does not mean that the
generated geometry was intrinsically better for all voxel detectors;
the architectures, initial checkpoints, and training procedures differ.

\begin{table}[t]\centering\setlength{\tabcolsep}{5pt}
\begin{tabular}{llrrrr}\toprule
Input & Weights & Easy & Moderate & Hard & $\Delta$ Mod. \\
\midrule\tableinput{tables/kitti_centerpoint_B.tex}\bottomrule
\end{tabular}
\caption{CenterPoint Line~B Car 3D AP$_{R40}$, 3,769 frames per row.
Frozen and adapted arms use their corresponding sparse-baseline control.}
\label{tab:cp-b}\end{table}

The three-epoch CenterPoint findings were not confined to Cars. For Line~A,
observed-first adaptation reached Pedestrian and Cyclist Moderate AP
of 48.42 and 57.70, below the original baseline's 50.65 and 64.61.
For Line~B the corresponding values were 44.06 and 33.56, compared
with 47.44 and 42.07 for the adapted sparse baseline.
Figure~\ref{fig:k-classes} displays these class-specific comparisons.
The three-epoch observed-first adaptation therefore did not exceed the
appropriate reference for any of these three Moderate 3D metrics.

\thesisfigure{k_classes}{Class-specific adaptation outcome}{CenterPoint
Moderate 3D AP$_{R40}$ for the three-epoch observed-first adapted representation
and its reference. Line~A uses the frozen original baseline; Line~B uses
the adapted sparse baseline. All values use 3,769 frames.}{fig:k-classes}

\FloatBarrier
\subsection{Class, Difficulty, and Projection Results}
\label{sec:kitti-full}

The Car 3D results are already complete in
Tables~\ref{tab:pr-a}--\ref{tab:cp-b} and are not repeated here.
Tables~\ref{tab:pr-bev-full} and~\ref{tab:pr-bbox-full} give PointRCNN
bird's-eye-view (BEV) and image-box AP, while
Table~\ref{tab:cp-noncar-full} supplies the remaining CenterPoint
Pedestrian and Cyclist 3D results and Table~\ref{tab:cp-bev-full}
its BEV results. Values are AP$_{R40}$ percentages over all 3,769
frames; F and A denote frozen and adapted detector weights.
Direct and observed-first refer to the PU-GCN input construction.
CenterPoint image-box AP was not exported.

\begin{table}[!t]\centering
\begin{tabular}{llrrr}\toprule
Input / weights & Class & Easy & Moderate & Hard\\\midrule
A / Baseline / F & Car & 95.87 & 88.86 & 86.43 \\
A / Direct / F & Car & 91.65 & 71.45 & 68.80 \\
A / Direct / A & Car & 92.28 & 80.52 & 76.09 \\
A / Observed / F & Car & 92.18 & 73.79 & 70.89 \\
A / Observed / A & Car & 94.00 & 81.03 & 78.32 \\
B / Baseline / F & Car & 91.53 & 74.92 & 71.90 \\
B / Baseline / A & Car & 91.75 & 77.97 & 73.67 \\
B / Direct / F & Car & 59.93 & 38.74 & 32.48 \\
B / Direct / A & Car & 82.49 & 58.20 & 51.48 \\
B / Observed / F & Car & 65.65 & 43.86 & 38.93 \\
B / Observed / A & Car & 87.48 & 67.30 & 60.73 \\
\bottomrule\end{tabular}
\caption{PointRCNN BEV AP$_{R40}$ on the final validation set.}\label{tab:pr-bev-full}
\end{table}

\begin{table}[!t]\centering
\begin{tabular}{llrrr}\toprule
Input / weights & Class & Easy & Moderate & Hard\\\midrule
A / Baseline / F & Car & 99.19 & 92.47 & 89.96 \\
A / Direct / F & Car & 95.52 & 74.57 & 71.80 \\
A / Direct / A & Car & 95.57 & 84.01 & 79.40 \\
A / Observed / F & Car & 95.78 & 77.21 & 74.37 \\
A / Observed / A & Car & 97.32 & 86.18 & 81.70 \\
B / Baseline / F & Car & 95.66 & 80.04 & 75.59 \\
B / Baseline / A & Car & 95.31 & 81.52 & 78.86 \\
B / Direct / F & Car & 63.38 & 39.83 & 34.99 \\
B / Direct / A & Car & 88.58 & 61.78 & 54.84 \\
B / Observed / F & Car & 69.00 & 46.46 & 39.98 \\
B / Observed / A & Car & 91.78 & 71.12 & 64.30 \\
\bottomrule\end{tabular}
\caption{PointRCNN Image-Box AP$_{R40}$ on the final validation set.}\label{tab:pr-bbox-full}
\end{table}

\begin{table}[!t]\centering
\begin{tabular}{llrrr}\toprule
Input / weights & Class & Easy & Moderate & Hard\\\midrule
A / Baseline / F & Pedestrian & 54.04 & 50.65 & 46.25 \\
A / Baseline / F & Cyclist & 79.24 & 64.61 & 60.99 \\
A / Direct / F & Pedestrian & 49.07 & 45.08 & 41.06 \\
A / Direct / F & Cyclist & 68.40 & 42.83 & 40.32 \\
A / Observed / F & Pedestrian & 49.61 & 45.98 & 42.02 \\
A / Observed / F & Cyclist & 73.38 & 49.45 & 46.49 \\
A / Observed / A & Pedestrian & 51.96 & 48.42 & 44.17 \\
A / Observed / A & Cyclist & 77.96 & 57.70 & 54.06 \\
B / Baseline / F & Pedestrian & 27.20 & 24.57 & 21.67 \\
B / Baseline / F & Cyclist & 25.81 & 15.46 & 14.49 \\
B / Baseline / A & Pedestrian & 51.70 & 47.44 & 42.08 \\
B / Baseline / A & Cyclist & 63.17 & 42.07 & 39.46 \\
B / Direct / F & Pedestrian & 28.88 & 24.89 & 21.53 \\
B / Direct / F & Cyclist & 27.12 & 15.28 & 14.36 \\
B / Observed / F & Pedestrian & 32.55 & 29.03 & 24.99 \\
B / Observed / F & Cyclist & 34.28 & 20.02 & 18.82 \\
B / Observed / A & Pedestrian & 48.88 & 44.06 & 38.98 \\
B / Observed / A & Cyclist & 55.80 & 33.56 & 31.59 \\
\bottomrule\end{tabular}
\caption{CenterPoint Pedestrian and Cyclist 3D AP$_{R40}$ on the final validation set.}\label{tab:cp-noncar-full}
\end{table}

\begin{table}[!t]\centering
\begin{tabular}{llrrr}\toprule
Input / weights & Class & Easy & Moderate & Hard\\\midrule
A / Baseline / F & Car & 92.18 & 88.43 & 87.03 \\
A / Baseline / F & Pedestrian & 58.98 & 55.74 & 51.62 \\
A / Baseline / F & Cyclist & 82.38 & 68.70 & 65.26 \\
A / Direct / F & Car & 88.99 & 70.68 & 68.96 \\
A / Direct / F & Pedestrian & 53.09 & 49.38 & 45.80 \\
A / Direct / F & Cyclist & 72.31 & 46.71 & 44.09 \\
A / Observed / F & Car & 89.77 & 72.83 & 71.23 \\
A / Observed / F & Pedestrian & 54.17 & 50.82 & 46.94 \\
A / Observed / F & Cyclist & 76.27 & 52.44 & 49.51 \\
A / Observed / A & Car & 91.85 & 84.93 & 83.68 \\
A / Observed / A & Pedestrian & 55.58 & 52.61 & 48.34 \\
A / Observed / A & Cyclist & 81.11 & 60.87 & 57.30 \\
B / Baseline / F & Car & 89.12 & 77.31 & 73.16 \\
B / Baseline / F & Pedestrian & 34.49 & 31.55 & 28.01 \\
B / Baseline / F & Cyclist & 28.27 & 17.13 & 16.33 \\
B / Baseline / A & Car & 90.17 & 80.20 & 76.79 \\
B / Baseline / A & Pedestrian & 58.01 & 53.82 & 48.39 \\
B / Baseline / A & Cyclist & 66.05 & 45.01 & 42.69 \\
B / Direct / F & Car & 74.51 & 47.53 & 42.86 \\
B / Direct / F & Pedestrian & 33.53 & 29.56 & 25.75 \\
B / Direct / F & Cyclist & 29.98 & 17.00 & 16.42 \\
B / Observed / F & Car & 80.99 & 55.54 & 50.94 \\
B / Observed / F & Pedestrian & 37.72 & 33.62 & 29.55 \\
B / Observed / F & Cyclist & 36.27 & 21.81 & 20.86 \\
B / Observed / A & Car & 88.20 & 73.67 & 69.72 \\
B / Observed / A & Pedestrian & 54.66 & 49.66 & 44.12 \\
B / Observed / A & Cyclist & 57.93 & 35.62 & 33.93 \\
\bottomrule\end{tabular}
\caption{CenterPoint BEV AP$_{R40}$ on the final validation set.}\label{tab:cp-bev-full}
\end{table}


\subsection{PU-GCN Inputs and Detector Retraining}
\label{sec:adaptation-evidence}

The training work updated detector weights while keeping PU-GCN fixed.
It progressed from a small adaptation screen to the three-epoch,
full-validation matrix and then to extended observed-first training.
The source training split contained 3,712 frames; the later PointRCNN
loader retained 3,265 after filtering. All main comparisons below use
3,769 validation frames. The earlier 256-frame pilot remains a separate
development experiment rather than part of the full-validation result.

The three-epoch full-validation experiment followed a smaller adaptation screen, which should
remain visible as completed work. It used 64 selected training frames
(63 effective after filtering), seed 20260823, three RPN epochs and
three RCNN epochs. Six input conditions were evaluated before and after
adaptation on the same 256 validation frames; all twelve evaluations
completed. Table~\ref{tab:k-screen} reports their Moderate results under
the explicitly recorded legacy Python evaluator. Change is adapted
minus frozen AP; Gap is adapted AP minus the matching adapted baseline,
both in AP points. These values are not points on an R40 learning curve.
The complete Easy/Moderate/Hard triplets are retained in the archived
experiment outputs.

\begin{table}[!t]\centering\setlength{\tabcolsep}{4pt}
\begin{tabular}{llrrrr}\toprule
Line & Input & Frozen & Adapted & Change & Gap\\\midrule
A & Original & 79.19 & 78.18 & $-1.01$ & 0.00\\
A & PU-GCN & 60.27 & 64.44 & $+4.17$ & $-13.74$\\
A & PDANS & 68.58 & 67.02 & $-1.57$ & $-11.17$\\
B & Sparse & 66.62 & 62.91 & $-3.70$ & 0.00\\
B & PU-GCN & 32.92 & 37.52 & $+4.59$ & $-25.40$\\
B & PDANS & 43.01 & 47.43 & $+4.43$ & $-15.48$\\\bottomrule
\end{tabular}
\caption{PointRCNN adaptation screen on 256 frames: Car Moderate 3D AP
at IoU 0.70 under the legacy evaluator.}\label{tab:k-screen}
\end{table}

The screen recovered more than four AP points for both PU-GCN inputs
and for sparse PDANS, but not for dense PDANS. Both adapted baselines
also deteriorated. Thus a positive before/after change alone would
overstate the benefit: all adapted generated inputs remained below
their paired adapted baseline. The later stage expanded training to
3,712 frames, initially evaluated a 256-frame pilot, and ultimately
completed the 3,769-frame, 20-arm matrix. Increasing training coverage,
changing the integration configuration and changing the evaluation
scope prevent these stages from being interpreted as a single
controlled sample-size ablation.


The training losses decreased. For PointRCNN Line~A direct input, median
RPN loss changed from 1.624055 to 1.347459 and RCNN loss from 1.043027
to 0.959636 between epochs~1 and~3. CenterPoint loss changed from
2.33 to 2.11 for Line~A observed-first, from 3.69 to 3.10 for Line~B
observed-first, and from 3.00 to 2.65 for the sparse baseline.
These were successful optimisation records, but the original matrix retained
only epoch~3 checkpoints and recorded no epoch~1/2 validation AP. That
fixed-duration study alone did not establish a validation plateau.

\phantomsection\label{sec:kitti-extended}
The extended observed-first experiments improved all four selected
validation scores beyond their three-epoch counterparts, but none
reached its line-specific reference
(Table~\ref{tab:k-extended-comparison}). The comparison separates
detector adaptation on an unchanged upsampled input from comparison
with the original observation. Thus the frozen observed-first rows
measure the starting point for adaptation, whereas the original-$N$
and sparse-$M$ rows are distinct input controls. The original-$N$ row
also serves as the complete-observation reference for Line~B.

Frozen denotes official detector weights. Extended denotes the
validation-selected best checkpoint from the schedules in
Table~\ref{tab:k-extended-stages}, with the PU-GCN checkpoint unchanged.
Baseline training budgets were not extended.

\begin{table}[!t]\centering\setlength{\tabcolsep}{4pt}
\begin{tabular}{lL{0.43\linewidth}rr}\toprule
Line & Input / detector weights & PointRCNN & CenterPoint\\\midrule
\tableinput{tables/kitti_extended_comparison.tex}
\bottomrule\end{tabular}
\caption{Detector-adaptation comparison: Car Moderate 3D AP$_{R40}$ on
3,769 validation frames. Bold marks extended-schedule results,
not the highest scores across input controls.}
\label{tab:k-extended-comparison}\end{table}

Relative to the same observed-first input with frozen detector weights,
the extended PointRCNN results gained 7.32 AP points in Line~A and
21.52 in Line~B. The corresponding CenterPoint gains were 12.76 and
18.81 points. The additional gains over the earlier three-epoch
adaptation were smaller: 0.58 and 1.83 for PointRCNN, and 1.26 and
1.24 for CenterPoint. Figure~\ref{fig:k-convergence-gains} separates
these recovery gains from the remaining baseline deficits, so a large
improvement from a degraded starting point is not mistaken for an
advantage over the original input.

\thesisfigure{k_convergence_gains}{Recovery and remaining deficits after extended adaptation}
{Car Moderate 3D AP$_{R40}$ differences on 3,769 validation frames.
(a) Extended-best minus frozen observed-first and minus three-epoch
observed-first results. (b) Extended-best minus the original-$N$ frozen
baseline and minus the line-specific reference. The latter is the same
original baseline in Line~A and the three-epoch adapted sparse-$M$
baseline in Line~B; the two markers therefore coincide in Line~A.
Each point comes from one completed schedule and its recorded comparison,
not an average across independent seeds.}{fig:k-convergence-gains}

For Line~A, the remaining original-baseline deficits were 10.36 AP points
for PointRCNN and 3.31 for CenterPoint. In Line~B, the extended results
remained 11.28 and 5.15 points below the adapted sparse controls.
They were also below the \emph{unadapted} sparse controls, by 8.53 and
1.70 points, respectively. Compared with the complete original
observation, the Line~B deficits were larger still, at 24.90 and
16.38 points. The recovery therefore does not depend on choosing only
one unfavourable starting comparator.

The stage summaries in Table~\ref{tab:k-extended-stages} establish the
scope of the plateau assessment. PointRCNN selected RPN epoch~14 in
both lines; the associated RCNN selections were epoch~2 in Line~A and
epoch~19 in the 24-epoch Line~B schedule. The four final PointRCNN
trajectories contain 78 epoch evaluations in total. CenterPoint selected
epoch~12 in each of its two 12-epoch schedules, adding 24 evaluations.
These 102 checkpoint evaluations are repeated measurements of the
completed training trajectories, not 102 independent training runs.
PointRCNN RPN-stage AP belongs to its stage-specific detection
evaluation and is not an additional final RCNN result.

\begin{table}[!t]\centering\setlength{\tabcolsep}{4pt}
\begin{tabular}{lllrrrr}\toprule
Detector & Line & Stage & \shortstack{Total\\epochs} &
\shortstack{Best\\epoch} & \shortstack{Best\\AP} &
\shortstack{Final\\count}\\\midrule
\tableinput{tables/kitti_extended_stages.tex}
\bottomrule\end{tabular}
\caption{Completed schedules and selected checkpoints. Final count denotes
consecutive terminal epochs without improvement exceeding the running
reference by 0.2 AP points (Section~\ref{sec:fc-4-11}).}
\label{tab:k-extended-stages}\end{table}

All six terminal counts in Table~\ref{tab:k-extended-stages} meet the
threshold of three. A numerical maximum at the last epoch does not
contradict this threshold-based plateau rule. CenterPoint Line~A improved from the last
threshold-reference score of 75.7998 at epoch~8 to 75.9661 at epoch~12;
Line~B improved from 62.8977 at epoch~9 to 62.9007 at epoch~12.
These increases of 0.1663 and 0.0030 AP points did not reset the counter.
For PointRCNN Line~B, Figure~\ref{fig:k-convergence-terminal} shows
the six recorded terminal RCNN evaluations. The selected epoch~19
scored 57.0485, whereas epoch~24 scored 55.5902 and all five intervening
or final evaluations remained below the selected value. The last
checkpoint was therefore not the selected checkpoint.

\thesisfigure{k_convergence_terminal}{Terminal validation after the selected PointRCNN checkpoint}
{PointRCNN Line~B Car Moderate 3D AP$_{R40}$ for epochs~19--24
of the completed 24-epoch RCNN schedule, with RPN epoch~14 fixed.
Blue marks the selected RCNN epoch; grey marks subsequent epochs.
The panel shows the complete terminal six-epoch window tabulated in
the experiment record, rather than the full training trajectory.
Every point evaluates all 3,769 validation frames; the connecting
segments join adjacent recorded epochs and do not estimate missing
earlier values.}{fig:k-convergence-terminal}

This extension supplies the validation-progress evidence missing from
the initial three-epoch study. It narrows the explanation for the
remaining deficit: stopping after three epochs alone does not account
for all of the loss under the tested schedules. It does not establish
a global optimum, exclude improvement with another optimiser or input
design, or remove the unequal-budget and validation-selection limits.
The earlier all-class tables retain their original three-epoch
checkpoints; their Pedestrian and Cyclist values are not reassigned
to these Car-selected models.


\FloatBarrier
\section{KITTI Integration and Mechanism Diagnostics}
\label{sec:kitti-diagnostics}

\subsection{Earlier Multi-Method Evidence}

The preceding multi-method experiment is informative about integration,
but it is a distinct experimental generation. Its CenterPoint evaluations
covered all 3,769 validation frames under AP$_{R40}$.
Figure~\ref{fig:k-legacy} retains all three Moderate class results.
PDANS gave the highest Car AP among the four generated-input pipelines
in both lines, yet remained below the corresponding Car baseline.
The class-specific exception matters: on Line~B, PDANS Pedestrian AP
was 28.18 against the sparse baseline's 24.57, and its Cyclist AP
was 15.52 against 15.46. These results rule out the stronger claim
that every class always deteriorated.

\thesisfigure{k_legacy}{Earlier CenterPoint multi-method comparison}{
Recorded full-validation Moderate 3D AP$_{R40}$ from the earlier
integration generation. PU-Net old contains a known normalisation error;
PU-EdgeFormer compat uses the lab compatibility path. These labels
identify implemented pipelines and do not establish a fair ranking of
the original architectures.}{fig:k-legacy}

Two implementation qualifications constrain that matrix. The old PU-Net
adapter had a confirmed normalisation error, so its very low AP cannot
be attributed to the PU-Net architecture. The lab PU-EdgeFormer path
reused compatible PU-GCN operators and checkpoint inference components;
it is not the independently verified native EdgeTransformer execution
used for the HPC ModelNet40 branch. The historical PointRCNN E2/E3
results can be identified more precisely: their archived runner uses
C++ AP$_{R40}$. They are therefore included below with their own paired
baselines, while the legacy adaptation screen retains its different
metric label.

\subsubsection{PointRCNN E2 and E3 Full-Validation Controls}
\label{sec:k-e2-e3}

The earlier E2 protocol evaluated a deterministic 16,384-point input
constructed after field-of-view and range filtering, using 0.1~m voxel
representatives and depth-stratified quotas. Its four depth intervals
were 0--20, 20--40, 40--60 and 60--70.4~m; largest-remainder allocation
assigned quotas from candidate counts without labels or AP-based
selection. E3 tested a 32,768-point, observed-first input. Both completed
ten conditions on all 3,769 validation frames. The primary CSV and
runner identify these as C++ AP$_{R40}$ results.

Figure~\ref{fig:k-e2-e3} and Tables~\ref{tab:k-e2-3d}--\ref{tab:k-e3}
retain this independent comparison. E2's original and sparse Car
Moderate baselines were 82.26 and 65.75. Among generated inputs, PDANS
was highest in both lines at 67.22 and 44.08. These figures describe
the integrated pipelines, including the old PU-Net error and
PU-EdgeFormer compatibility path; they do not rank faithful original
implementations.

\thesisfigure{k_e2_e3}{PointRCNN input-budget sensitivity}{Car Moderate
3D AP$_{R40}$ for the historical E2 (16,384-point) and E3 (32,768-point,
observed-first) inputs, each evaluated on 3,769 frames. Markers share a
0--100 axis. The change combines input budget and selection policy;
it is not a pure point-count effect. Baselines belong to these protocols,
not the later PU-GCN adaptation matrix.}{fig:k-e2-e3}

\begin{table}[!t]\centering\setlength{\tabcolsep}{5pt}
\begin{tabular}{llrrr}\toprule
Line & Input & Easy & Moderate & Hard\\\midrule

A & Baseline & 92.27 & 82.26 & 77.95\\

A & PDANS & 84.61 & 67.22 & 62.40\\

A & PU-GCN & 80.17 & 58.28 & 53.47\\

A & PU-EF compat & 65.68 & 44.96 & 40.19\\

A & PU-Net old & 11.76 & 8.38 & 7.44\\

\midrule

B & Baseline & 85.18 & 65.75 & 61.38\\

B & PDANS & 65.01 & 44.08 & 39.17\\

B & PU-GCN & 45.93 & 28.68 & 24.13\\

B & PU-EF compat & 33.76 & 20.57 & 17.62\\

B & PU-Net old & 12.41 & 8.78 & 7.76\\

\bottomrule\end{tabular}
\caption{Historical PointRCNN E2 3D AP$_{R40}$ for Car. Each row covers all 3,769 frames; the integrated input protocol is shared within E2. Values are percentages.}\label{tab:k-e2-3d}\end{table}

\begin{table}[!t]\centering\setlength{\tabcolsep}{5pt}
\begin{tabular}{llrrr}\toprule
Line & Input & Easy & Moderate & Hard\\\midrule

A & Baseline & 95.98 & 88.98 & 86.59\\

A & PDANS & 90.00 & 78.02 & 73.21\\

A & PU-GCN & 86.93 & 66.79 & 62.02\\

A & PU-EF compat & 75.97 & 53.83 & 48.94\\

A & PU-Net old & 28.95 & 18.74 & 16.85\\

\midrule

B & Baseline & 91.04 & 76.40 & 72.03\\

B & PDANS & 76.63 & 54.57 & 48.14\\

B & PU-GCN & 57.25 & 36.19 & 31.43\\

B & PU-EF compat & 43.65 & 27.74 & 23.38\\

B & PU-Net old & 23.77 & 15.52 & 13.62\\

\bottomrule\end{tabular}
\caption{Historical PointRCNN E2 BEV AP$_{R40}$ for Car. Each row covers all 3,769 frames; the integrated input protocol is shared within E2. Values are percentages.}\label{tab:k-e2-bev}\end{table}

\begin{table}[!t]\centering\setlength{\tabcolsep}{5pt}
\begin{tabular}{llrrr}\toprule
Line & Input & Easy & Moderate & Hard\\\midrule

A & Baseline & 99.32 & 94.08 & 90.10\\

A & PDANS & 93.32 & 81.75 & 76.85\\

A & PU-GCN & 92.44 & 71.77 & 67.07\\

A & PU-EF compat & 81.95 & 58.80 & 52.30\\

A & PU-Net old & 38.46 & 23.89 & 20.69\\

\midrule

B & Baseline & 95.24 & 79.95 & 75.58\\

B & PDANS & 83.75 & 60.33 & 53.79\\

B & PU-GCN & 60.61 & 38.61 & 32.40\\

B & PU-EF compat & 52.75 & 33.49 & 28.71\\

B & PU-Net old & 33.62 & 20.51 & 18.40\\

\bottomrule\end{tabular}
\caption{Historical PointRCNN E2 image-box AP$_{R40}$ for Car. Each row covers all 3,769 frames; the integrated input protocol is shared within E2. Values are percentages.}\label{tab:k-e2-bbox}\end{table}

\begin{table}[!t]\centering\setlength{\tabcolsep}{4pt}
\begin{tabular}{llrrrr}\toprule
Line & Input & Easy & Moderate & Hard & $\Delta$ Mod.\\\midrule

A & Baseline & 92.03 & 80.38 & 77.52 & -1.88\\

A & PDANS & 86.27 & 70.06 & 65.27 & +2.84\\

A & PU-GCN & 85.19 & 64.95 & 58.66 & +6.68\\

A & PU-EF compat & 71.54 & 51.63 & 45.65 & +6.67\\

A & PU-Net old & 27.92 & 18.93 & 17.88 & +10.55\\

\midrule

B & Baseline & 84.59 & 64.14 & 57.50 & -1.61\\

B & PDANS & 64.92 & 44.04 & 39.16 & -0.04\\

B & PU-GCN & 45.09 & 28.28 & 23.83 & -0.40\\

B & PU-EF compat & 32.50 & 20.03 & 16.20 & -0.54\\

B & PU-Net old & 10.99 & 7.48 & 7.19 & -1.31\\

\bottomrule\end{tabular}
\caption{Historical PointRCNN E3 Car 3D AP$_{R40}$. $\Delta$ Mod. is E3
minus E2 for the same input name, in AP points. All ten rows have 3,769
prediction files.}\label{tab:k-e3}\end{table}

E3 raised dense-input PDANS from 67.22 to 70.06 and PU-GCN from
58.28 to 64.95, while the original baseline fell from 82.26 to 80.38.
The stronger generated-input recovery therefore coexisted with a
changed baseline, and all Line~A generated inputs remained below it.
For Line~B, PDANS was approximately unchanged (44.08 to 44.04),
PU-GCN declined slightly (28.68 to 28.28), and the sparse baseline
fell from 65.75 to 64.14. Enlarging the budget with observed-first
selection did not remove the downstream deficit. Because both budget
and point selection changed, these results establish sensitivity of
the complete input adapter rather than a causal effect of cardinality
alone.

\FloatBarrier


\subsection{Patch Locality and Normalisation}
\label{sec:patch-diagnostics}

The original patch extractor could combine distant structures because
it partitioned coarse spatial bins by row order. In the audited Line~B
patches, the median XY diagonal was 30.46~m and its 90th percentile
was 124.10~m. Such supports do not resemble the local surface patches
used to train object upsamplers \cite{yu2018punet,qian2021pugcn}. A first local replacement reduced the
spatial extent but introduced 59.4\% duplicate rows. Locality and
unique support therefore had to be checked separately.

The first local PDANS replacement produced opposite detector responses
(Figure~\ref{fig:k-patch-response}). On the same 256 frames, PointRCNN
Car Moderate 3D AP fell from 57.18 to 43.70, while CenterPoint Car
rose from 61.12 to 77.42. The saved paired report also records
CenterPoint Pedestrian and Cyclist increases. Its locality rule guaranteed
only 256 unique supporting points for a 2,048-row patch; repeated filling
and overlapping patches concentrated the generated output. This was a
development configuration, distinct from the later unique-2,048-support
PU-GCN pipeline.

\thesisfigure{k_patch_response}{Paired PDANS patch responses}{Paired
Moderate 3D AP from the archived 256-frame PDANS patch experiment. The
source report labels these values AP$_{R40}$; they retain that report's
protocol and are not inserted into the final adaptation matrix. Each
segment connects the old and first local patch result for one detector
and class. The CenterPoint values use the original evaluation logs named
in the paired report, which differ slightly from an older summary CSV.}{fig:k-patch-response}

The retained input audit supports a concentration explanation. Across 20
frames, the median fraction of PointRCNN input coordinates unique at
1~mm fell from 99.90\% to 57.70\%; occupied diagnostic voxels fell from
11,584.5 to 2,124.5, and the share of points in the densest 10\% of voxels
rose from 27.23\% to 87.85\%. CenterPoint preserved the original measured
points before voxelisation, whereas PointRCNN sampled 16,384 input rows.
The Car matching audit changed from 539 to 402 matches for PointRCNN
and from 621 to 748 for CenterPoint. These counts corroborate the
opposite observed changes, but do not replace AP or isolate a single
causal operation: locality, duplication and effective support changed
together.

A 256-frame PU-Net ablation isolated normalisation and patch construction
as two factors (Figure~\ref{fig:k-norm}). Fixing normalisation with the
old patch improved PointRCNN Moderate AP from 7.98 to 18.27 and
CenterPoint from 15.98 to 41.22. Fixing both factors raised the values
to 43.54 and 46.20. The paired baselines were 81.30 and 79.84.
Both faults affected performance, but correcting them was insufficient
to recover the baselines. Because this was a diagnostic subset and an
earlier evaluator configuration, its AP values are not inserted into
the final full-validation ranking.

\thesisfigure{k_normalization}{Normalisation and patch ablation}{Recorded
256-frame PU-Net diagnostic: wrong or fixed normalisation crossed with
old or local patches. Values are compared only within this experiment.
The historical PointRCNN metric is not promoted to the final AP$_{R40}$
series.}{fig:k-norm}

\subsection{The Effective Detector Input}
\label{sec:budget-results}

The input-budget audits show why exact file cardinality is insufficient
to describe the detector's observation. PointRCNN samples a fixed 16,384
points at its entrance. CenterPoint limits points per voxel and the number
of retained voxels. In a 32-frame Line~A audit, median occupied-voxel
counts were 14,944 for Original, 36,913 for PDANS, 45,108 for PU-GCN,
46,182 for the PU-EdgeFormer compatibility path, and 55,384 for old
PU-Net. The 40,000-voxel cap was reached in 0, 6, 29, 29, and 32
frames, respectively (Figure~\ref{fig:k-voxels}).

\thesisfigure{k_voxels}{Voxel-budget audit}{Median occupied voxels and
numbers of cap-triggering frames in the archived 32-frame Line~A audit.
The dashed line is the 40,000-voxel evaluation cap. These are diagnostic
counts from the earlier input generation, not statistics of the final
3,769-frame adaptation matrix.}{fig:k-voxels}

Line~B did not trigger that cap in the audit, yet its detection losses
were large. Voxel truncation is therefore a relevant Line~A mechanism,
but cannot be the only explanation. Similarly, a separate full-input
PointRCNN experiment that removed fixed entrance resampling reduced even
the original-input Moderate result from approximately 81.95 to 74.91
within its own protocol. Simply presenting more points changed the input
distribution expected by the pretrained detector.

The region-split pilot tested a further alternative. It processed 2,067
shared core regions with a 3~m halo, preserved all core points, and kept
individual detector inputs within 16,384 points. Under that three-class
PointRCNN configuration, baseline Car/Pedestrian/Cyclist Moderate
AP$_{R40}$ was 77.82/58.21/78.29; the closest tested PDANS branch
reached 68.87/58.07/71.58. Avoiding core-point truncation did not
produce a consistent gain. This diagnostic uses a separate detector
implementation and inference protocol, so it cannot quantify the effect
of truncation in the final native PointRCNN experiment.

\subsection{Selection, Sampling Variance, and Negative Controls}
A retained 256-frame PDANS experiment tested a small generated-point
dose of 2.5\% against a matched observed-point fill control.
Car Moderate 3D AP was 82.6611 for the baseline, 84.4160 for the
observed-fill control, and 85.0953 for PDANS. The apparent 2.4342-point
gain over baseline therefore consisted of a 1.7548-point control
improvement and only a 0.6793-point additional difference from the
generated representation. This attribution prevents the whole gain
from being assigned to generated geometry.

Table~\ref{tab:k-dose-control} retains the other metrics from the same
comparison. The sign was not uniform across difficulty or projection:
the PDANS variant reduced Hard 3D AP and Moderate image-box AP.
This small-dose result was a diagnostic candidate rather than evidence
of an established full-validation gain. Values retain the historical
evaluator's metric definitions and are not pooled with the final
adaptation matrix. AOS denotes KITTI's average orientation similarity
\cite{geiger2012kitti}.

\begin{table}[!t]\centering\setlength{\tabcolsep}{4pt}
\begin{tabular}{llrrrr}\toprule
Metric & Difficulty & Baseline & Obs.\ fill & PDANS & $\Delta$ vs fill\\\midrule
BBox & Easy & 99.37 & 99.49 & 99.22 & -0.27 \\
BBox & Moderate & 94.12 & 93.93 & 92.31 & -1.62 \\
BBox & Hard & 89.59 & 89.48 & 89.53 & +0.05 \\
BEV & Easy & 96.45 & 96.53 & 96.43 & -0.10 \\
BEV & Moderate & 90.91 & 90.83 & 91.08 & +0.25 \\
BEV & Hard & 88.44 & 86.36 & 86.37 & +0.00 \\
3D & Easy & 95.09 & 94.80 & 95.44 & +0.64 \\
3D & Moderate & 82.66 & 84.42 & 85.10 & +0.68 \\
3D & Hard & 79.93 & 79.82 & 78.63 & -1.20 \\
AOS & Easy & 99.36 & 99.48 & 99.22 & -0.27 \\
AOS & Moderate & 94.01 & 93.80 & 92.19 & -1.61 \\
AOS & Hard & 89.44 & 89.33 & 89.34 & +0.02 \\
\bottomrule\end{tabular}
\caption{Low-dose comparison on 256 frames. The last column is PDANS
minus observed-fill control, in percentage points.}
\label{tab:k-dose-control}\end{table}

The executed detector-aware PDANS sequence tested whether generated
points could be restricted to more useful regions. The V4 measured-voxel
anchor reduced the median number of newly activated voxels to zero and
gave a 0.31-point Car gain in its PointRCNN pilot. CenterPoint Car
was approximately unchanged, while Pedestrian and Cyclist still declined.
The V5 class-adaptive gate nearly preserved all CenterPoint pilot metrics,
but its independent 64-frame holdout changed Car Moderate AP from
74.1881 to 74.0512. These tests support controlled occupancy as an
integration factor; they did not establish a repeatable cross-detector gain.

A 16-seed evaluator probe further qualified a positive single-run result.
With input, checkpoint, split, and merge rules fixed, changing only the
sampling seed gave Car Moderate scores (mean $\pm$ standard deviation across seeds) of $80.998\pm1.198$ for
the baseline and $80.560\pm1.086$ for the V1 full variant. The paired
difference was $-0.438\pm1.769$ AP points, and 9 of 16 differences
were negative. The previously reported 1.2227-point gain was the
variant's most favourable seed in that probe. These are evaluator-seed
statistics for the diagnostic path, not training-seed uncertainty and
not an error bar applicable to the final full-validation matrix.

\FloatBarrier
\subsubsection{Frame-Subset Size and AP Stability}
\label{sec:k-subset-scale}

A separate completed probe varied which validation frames were scored,
while keeping saved detector predictions fixed. It drew 200 subsets
without replacement at each size of 20, 50, 100, 256 and 512 frames,
using seed 0 and paired frame identities for the original, sparse and
sparse-PDANS inputs. Car Moderate 3D AP$_{R40}$ was recomputed from the
saved predictions. This is variation from frame selection, distinct
from the preceding 16-seed detector-input sampling probe and from
retraining variability.

Figure~\ref{fig:k-subset-scale} and Table~\ref{tab:k-subset-scale} show
the standard deviation of each paired AP contrast. For the sparse
minus original contrast, it fell from 6.02 points at 20 frames to
1.26 at 512 frames. For PDANS minus sparse, the standard deviation was
4.78 at 50 frames and 2.27 at 256 frames. Quoting the former contrast's
variability as uncertainty of the latter would misstate the experiment.

\thesisfigure{k_subset_scale}{Frame-subset size and paired AP variability}{
Standard deviation across 200 random frame subsets per size, calculated
from fixed historical PointRCNN predictions. Each curve is a different
paired Car Moderate 3D AP$_{R40}$ contrast. These are observed
subsampling standard deviations, not confidence intervals for the final
matrix and not results of additional detector training.}{fig:k-subset-scale}

\begin{table}[!t]\centering\setlength{\tabcolsep}{5pt}
\begin{tabular}{rrrr}\toprule
Frames & Sparse $-$ original & PDANS $-$ original & PDANS $-$ sparse\\\midrule

20 & 6.02 & 8.60 & 7.89\\

50 & 3.71 & 5.21 & 4.78\\

100 & 2.87 & 3.84 & 3.37\\

256 & 1.75 & 2.57 & 2.27\\

512 & 1.26 & 1.57 & 1.54\\

\bottomrule\end{tabular}
\caption{Sample standard deviations of paired Car Moderate AP contrasts
in AP points. All three contrasts use the same sampled frames within
a draw.}\label{tab:k-subset-scale}\end{table}

At 20 frames, the original-input AP's empirical 5th--95th percentile
range was 51.96--89.95, despite fixed model predictions. The mean was
75.63 compared with 81.99 on the complete source set. Small-subset AP
therefore changed both variability and average value, reflecting the
composition and recall support of the selected frames. This is a
reason to distinguish case diagnosis from full-validation evaluation.
The archived probe also multiplies standard deviation by 1.96 under a
sensitivity label. That number is only a descriptive scale here; it is
not a formal minimum detectable effect established by a power analysis.
No repeated-training uncertainty or extrapolated full-set confidence
interval is inferred from this probe.

\FloatBarrier

\subsection{Completed Input and Selection Interventions}
\label{sec:kitti-interventions}
The completed interventions changed input capacity, patch support,
generated-point dose, and selection rules. Their separate baselines
and evaluation scopes are retained below. Differences between tables
are not interpreted as isolated method effects when the pipelines
differ.

\subsubsection{Removing Entrance Resampling}
The all-valid PointRCNN experiment evaluated the original cloud and
four generated Line~A inputs without standard entrance resampling.
Table~\ref{tab:k-all-valid} shows that the original-input Moderate
result itself fell from about 81.95 to 74.91. PDANS was the strongest
generated input at 63.49, but still trailed the all-valid original.
Removing a point cap changed the representation expected by the
detector rather than simply releasing a performance bottleneck.

\begin{table}[!t]\centering\setlength{\tabcolsep}{4pt}
\begin{tabular}{lrrr}\toprule
Input & Easy & Moderate & Hard\\\midrule
Original / standard & 92.26 & 81.95 & 77.75\\
Original / all-valid & 79.83 & 74.91 & 72.73\\
PDANS / all-valid & 83.59 & 63.49 & 54.34\\
PU-GCN / all-valid & 78.21 & 54.26 & 47.11\\
PU-EF / all-valid & 67.11 & 43.85 & 37.10\\
PU-Net old / all-valid & 12.23 & 8.82 & 8.17\\
\bottomrule\end{tabular}
\caption{Recorded full-validation all-valid-input experiment. Values are Car 3D AP from that historical protocol. PU-EF is the compatibility path and PU-Net old retains the faulty adapter.}\label{tab:k-all-valid}\end{table}

\subsubsection{Correcting Local Patch Support across Methods}
Table~\ref{tab:k-patch-pairs} retains all four paired surface-patch
pilots and their two batch baselines. Every method improved with
surface-c32. Fixed-normalisation PU-Net additionally reached 44.0985
with cover-kNN, against its batch baseline of 81.7441. Local support
therefore affected the recorded performance, but did not close the
baseline gap.

\begin{table}[!t]\centering\setlength{\tabcolsep}{4pt}
\begin{tabular}{llrrr}\toprule
Batch & Method & Baseline & Old patch & Surface\\\midrule
Batch 1 & PDANS & 81.56 & 57.99 & 68.47\\
Batch 1 & PU-GCN & 81.56 & 55.81 & 63.45\\
Batch 2 & PU-EF & 81.74 & 41.74 & 58.02\\
Batch 2 & PU-Net fixed & 81.74 & 16.81 & 22.93\\
\bottomrule\end{tabular}
\caption{PointRCNN 256-frame paired patch pilots, Car Moderate 3D AP. Surface denotes surface-c32; the two pilot batches retain separate baselines.}\label{tab:k-patch-pairs}\end{table}

\subsubsection{Region Splitting with Complete Core-Point Retention}
The region-split pilot retained every core point across 2,067 shared
regions with a 3~m halo and at most 16,384 points per detector input.
Table~\ref{tab:k-region-all} reports every method and class.
PDANS was closest to the baseline, but no generated input improved
the three principal class outcomes together.

\begin{table}[!t]\centering\setlength{\tabcolsep}{4pt}
\begin{tabular}{lrrr}\toprule
Input & Car & Pedestrian & Cyclist\\\midrule
Baseline & 77.82 & 58.21 & 78.29\\
PDANS & 68.87 & 58.07 & 71.58\\
PU-GCN & 61.51 & 55.47 & 57.48\\
PU-EF & 61.28 & 56.14 & 56.18\\
PU-Net fixed & 26.18 & 25.02 & 28.92\\
\bottomrule\end{tabular}
\caption{Region-split pilot on 256 frames: Moderate 3D AP$_{R40}$. Generated rows use surface-c32; PU-Net normalisation is fixed. This is a separate three-class PointRCNN implementation and inference protocol.}\label{tab:k-region-all}\end{table}

An independent 64-frame region-split V4 holdout improved
Car/Pedestrian/Cyclist from 75.7267/38.1070/15.0000 to
77.4777/45.8565/17.2222. These positive results remain part of the
completed evidence, bounded by the small holdout and its region-split
protocol. In native whole-frame OpenPCDet PointRCNN, V4 changed Car
by $+4.0651$ under deterministic workers=0 but $-2.2718$ under
workers=2, with other class effects also changing sign. The records
therefore do not establish improvement stable across those sampling
implementations.

\subsubsection{Dose and Selector Components}
The 256-frame PDANS Line~A dose sequence used 2.5\%, 5\%, 7.5\%,
and 10\% generated points. Its reported Moderate results were
approximately 85.10, 82.37, 82.25, and 82.42, against 82.66.
All tested Line~B doses remained below the sparse reference.
Full-validation g10 runs completed for all four methods:
Line~A PDANS/PU-GCN/PU-EF/PU-Net old reached approximately
79.951/79.478/76.433/72.371, and Line~B reached
59.048/55.954/52.415/45.700. The g25 and g50 runs did not complete
and are not treated as evaluated results.

The V1 component experiment compared density, fixed dose, matched-real,
nearest-neighbour, and disabled-component variants
(Table~\ref{tab:k-components}). The full variant had the highest
single-evaluation score, but the matched-real control also improved.
The later evaluator-seed probe described above shows why this one
favourable score did not establish a stable gain.

\begin{table}[!t]\centering\setlength{\tabcolsep}{4pt}
\begin{tabular}{lr}\toprule
Configuration & Moderate\\\midrule
Baseline & 81.18\\
Density-only & 79.86\\
Fixed 2.5\% PDANS & 81.44\\
Full & 82.40\\
Matched-real dynamic & 81.80\\
NN-only & 81.59\\
No-adaptive & 80.15\\
No-confidence & 81.64\\
No-sparse & 80.94\\
\bottomrule\end{tabular}
\caption{PointRCNN V1 selector/component pilot on 256 frames, Car Moderate 3D AP$_{R40}$. The separate 16-seed probe qualifies these single-evaluation scores.}\label{tab:k-components}\end{table}

\subsubsection{Detector-Aware Selection Versions}
V2 used final-box guidance, V3 used internal pre-RCNN or pre-NMS
proposals, V4 anchored measured voxels, and V5 introduced class-adaptive
gating. PointRCNN V2 voxel-only and proposal-only scores were
79.3805 and 75.7533; their combined score was 79.2269.
Internal proposals in V3 yielded 79.2933, so that replacement did
not remove the loss. Table~\ref{tab:k-versions} retains the detector
and class comparisons.

\begin{table}[!t]\centering\setlength{\tabcolsep}{4pt}
\begin{tabular}{lrrrr}\toprule
Version & PR Car & CP Car & CP Ped. & CP Cyc.\\\midrule
Baseline & 81.18 & 79.84 & 43.51 & 76.11\\
V2 full & 79.23 & 76.77 & 48.56 & 72.28\\
V3 full & 79.29 & 76.74 & 48.59 & 72.29\\
V4 & 81.49 & 79.82 & 43.07 & 75.77\\
\bottomrule\end{tabular}
\caption{Completed 256-frame PDANS versions: PointRCNN Car and CenterPoint Car/Pedestrian/Cyclist Moderate 3D AP. CenterPoint Pedestrian/Cyclist baselines are obtained from the recorded scores and baseline deltas.}\label{tab:k-versions}\end{table}

V4 reduced the median number of newly activated voxels to zero,
yet did not improve all classes. V5 changed development-set
CenterPoint Car/Pedestrian/Cyclist by
$-0.0184/+0.0004/+0.0000$ and reused the PointRCNN V4 result.
The completed independent 64-frame holdout did not retain a positive
Car effect: 74.1881 became 74.0512, while BEV fell from 86.3893
to 85.2910. Input control recovered performance during development;
the holdout limits the claim of a general benefit.

\subsubsection{Line B Selection and Small Negative Probes}
Table~\ref{tab:k-lineb-selectors} compares the completed c2048-r4,
E1, consensus, and random pilots. For PointRCNN, PDANS favoured E1,
whereas PU-GCN and fixed PU-Net had their highest listed score with
consensus. For CenterPoint, consensus reduced PDANS, PU-GCN, and
PU-EF relative to random selection but improved PU-Net.
No listed generated input reached its corresponding baseline.

\begin{table}[!t]\centering\setlength{\tabcolsep}{4pt}
\begin{tabular}{lrrrrr}\toprule
Protocol & Base & PDANS & PU-GCN & PU-EF & PU-Net\\\midrule
\shortstack[l]{PointRCNN\\c2048-r4} & 64.04 & 36.62 & 24.96 & 21.24 & 2.40\\
\shortstack[l]{PointRCNN\\E1} & 64.04 & 44.05 & 32.84 & 28.26 & 10.58\\
\shortstack[l]{PointRCNN\\consensus} & 63.42 & 39.25 & 33.16 & 22.44 & 18.88\\
\shortstack[l]{CenterPoint\\random} & 61.66 & 47.74 & 41.38 & 35.43 & 15.15\\
\shortstack[l]{CenterPoint\\consensus} & 61.66 & 43.61 & 37.31 & 27.33 & 19.80\\
\bottomrule\end{tabular}
\caption{Line~B 256-frame selection pilots, Car Moderate 3D AP. PU-EF is the compatibility implementation and PU-Net uses fixed normalisation. Protocol-specific baselines are not pooled.}\label{tab:k-lineb-selectors}\end{table}

A quality-selection probe used three frames without a voxel quota
and five frames with quota 6. Without a quota, reported quality
precision was 0.6821 versus 0.8155 for random selection, and unsupported
fraction was 0.5825 versus 0.4940. With quota 6, the comparisons
were 0.5965 versus 0.8024 and 0.6329 versus 0.5337.
These small geometric probes had no detector AP and were not expanded.
Their negative outcomes qualify the selection hypothesis, while saying
nothing directly about detector performance.

Together, these interventions tested capacity, local support, dose,
proposal guidance, measured-voxel anchoring, and selector components.
Several recovered performance within their own protocols, but no
consistent cross-protocol gain survived the available controls.
They motivated the later PU-GCN input construction and detector
adaptation, whose final full-validation results remain separate.

\FloatBarrier
\subsection{Object-Level Losses and Recoveries across the Validation Set}
\label{sec:kitti-transitions}

The historical lab audit retained individual object matches across the
3,769-frame validation split. This evidence connects the aggregate
degradation to concrete losses and recoveries. It belongs to the earlier
input-preserving integration generation, with the old PU-Net
normalisation defect and the PU-EdgeFormer compatibility path.
It is therefore analysed within that generation, separately from the
final PU-GCN adaptation matrix.

Each object was greedily matched to a same-class prediction by oriented
3D IoU. Car required 0.70, while Pedestrian and Cyclist required 0.50.
A lost object had a qualifying baseline match and none after upsampling;
a recovered object had the reverse transition. Maintained,
confidence-degraded, and localisation-degraded objects remained matched.
The audit does not reproduce the official difficulty filters, DontCare
treatment, or score-threshold integration. Its rates describe these
retained matches and are not AP or official difficulty-specific recall.

For PointRCNN Line~A, PU-GCN lost 3,597 baseline-matched Cars and recovered
210, giving a net change of $-3,387$ matched objects.
The loss rate was 32.65\% of 11,016 baseline matches.
For Line~B, it lost 4,829 and recovered 115: 55.37\% of the
8,721 baseline matches were lost. PDANS had a smaller loss in both
lines, but its recoveries also failed to offset its losses.
Figure~\ref{fig:k-transition-pr} makes both directions visible,
so occasional improvements are not concealed by a negative aggregate.

\thesisfigure{k_transitions_pointrcnn}{PointRCNN lost and recovered objects}{
Historical full-validation Car matching audit. Lost and recovered counts
use the same 14,385 Car opportunities within each comparison.
PU-Net* denotes the historical faulty adapter; PU-EF denotes the lab
compatibility path. These are object-event counts, not AP.}{fig:k-transition-pr}

Figure~\ref{fig:k-transition-cp} shows the corresponding CenterPoint counts. For Line~A, PU-GCN lost 2,932 Cars and recovered 340;
Line~B lost 3,505 and recovered 333. Its loss rates of 25.88\% and
37.43\% were lower than PointRCNN's in the corresponding historical
comparison, but both remained negative in net matches.
The two detectors thus share a qualitative failure pattern while
differing in its magnitude. This does not isolate architecture as the
cause: their checkpoints, preprocessing, and native inputs also differ.

\thesisfigure{k_transitions_centerpoint}{CenterPoint lost and recovered objects}{
Historical full-validation Car matching counts. The comparison retains
both loss and recovery events for every implemented pipeline.
Definitions and implementation qualifications match
Figure~\ref{fig:k-transition-pr}.}{fig:k-transition-cp}

Table~\ref{tab:k-transition-all} retains every class and method row
from this audit. PointRCNN is Car-only in this configuration; CenterPoint
also supplies Pedestrian and Cyclist. Different class thresholds and
baseline match counts are kept explicit. A recovery count is not
normalised by baseline TP: its natural denominator is baseline FN.
Loss and recovery percentages therefore cannot be subtracted directly.
The net count equals recovered minus lost.

\begin{longtable}{lllrrrr}
\caption{Complete retained transition summary. $B$ is baseline matched
count, $L$ is lost, $R$ is recovered, and $\Delta=R-L$.
PU-EF is the compatibility path and PU-Net* the faulty historical adapter.}
\label{tab:k-transition-all}\\
\toprule Line & Method & Class & $B$ & $L$ & $R$ & $\Delta$\\\midrule\endfirsthead
\toprule Line & Method & Class & $B$ & $L$ & $R$ & $\Delta$\\\midrule\endhead
\bottomrule\endfoot
\multicolumn{7}{l}{\textit{PointRCNN}}\\
A & PDANS & Car & 11,016 & 2,544 & 353 & -2,191 \\
A & PU-EF & Car & 11,016 & 5,279 & 143 & -5,136 \\
A & PU-GCN & Car & 11,016 & 3,597 & 210 & -3,387 \\
A & PU-Net* & Car & 11,016 & 8,909 & 39 & -8,870 \\
B & PDANS & Car & 8,721 & 3,091 & 317 & -2,774 \\
B & PU-EF & Car & 8,721 & 5,845 & 94 & -5,751 \\
B & PU-GCN & Car & 8,721 & 4,829 & 115 & -4,714 \\
B & PU-Net* & Car & 8,721 & 6,668 & 45 & -6,623 \\
\multicolumn{7}{l}{\textit{CenterPoint}}\\
A & PDANS & Car & 11,329 & 2,330 & 424 & -1,906 \\
A & PDANS & Cyclist & 595 & 113 & 58 & -55 \\
A & PDANS & Pedestrian & 1,568 & 277 & 117 & -160 \\
A & PU-EF & Car & 11,329 & 4,447 & 255 & -4,192 \\
A & PU-EF & Cyclist & 595 & 121 & 67 & -54 \\
A & PU-EF & Pedestrian & 1,568 & 538 & 112 & -426 \\
A & PU-GCN & Car & 11,329 & 2,932 & 340 & -2,592 \\
A & PU-GCN & Cyclist & 595 & 120 & 75 & -45 \\
A & PU-GCN & Pedestrian & 1,568 & 340 & 113 & -227 \\
A & PU-Net* & Car & 11,329 & 8,016 & 98 & -7,918 \\
A & PU-Net* & Cyclist & 595 & 268 & 34 & -234 \\
A & PU-Net* & Pedestrian & 1,568 & 755 & 107 & -648 \\
B & PDANS & Car & 9,364 & 2,587 & 493 & -2,094 \\
B & PDANS & Cyclist & 272 & 100 & 102 & +2 \\
B & PDANS & Pedestrian & 1,060 & 302 & 268 & -34 \\
B & PU-EF & Car & 9,364 & 5,175 & 227 & -4,948 \\
B & PU-EF & Cyclist & 272 & 123 & 77 & -46 \\
B & PU-EF & Pedestrian & 1,060 & 575 & 172 & -403 \\
B & PU-GCN & Car & 9,364 & 3,505 & 333 & -3,172 \\
B & PU-GCN & Cyclist & 272 & 81 & 94 & +13 \\
B & PU-GCN & Pedestrian & 1,060 & 372 & 203 & -169 \\
B & PU-Net* & Car & 9,364 & 6,274 & 148 & -6,126 \\
B & PU-Net* & Cyclist & 272 & 154 & 40 & -114 \\
B & PU-Net* & Pedestrian & 1,060 & 608 & 148 & -460 \\
\end{longtable}

\subsection{Where Matched Objects Are Lost}
\label{sec:kitti-distance}

Distance stratification asks whether the failure is concentrated in
sparser, more distant observations. The archived audit partitions
objects into planar ranges of 0--20, 20--40, and at least 40~m.
Figures~\ref{fig:k-distance-pr} and~\ref{fig:k-distance-cp} retain all
three bins and all four pipelines. Each point is lost baseline matches
divided by baseline matches in that same bin, not a fraction of all
objects or all detections.

For PointRCNN Line~A, PU-GCN lost 9.5\%, 44.6\%, and 81.6\% of baseline-matched Cars in these three bins, respectively.
For PointRCNN Line~B, PU-GCN lost 31.2\%, 86.8\%, and 95.4\% of baseline-matched Cars in these three bins, respectively.
For CenterPoint Line~A, PU-GCN lost 7.2\%, 33.0\%, and 69.2\% of baseline-matched Cars in these three bins, respectively.
For CenterPoint Line~B, PU-GCN lost 15.3\%, 58.8\%, and 84.3\% of baseline-matched Cars in these three bins, respectively.

The distance curves locate the degradation but do not establish
distance as its sole cause. Occlusion, object size, pose, visible
surfaces, and detector confidence also vary across range bins.
The conditional denominator additionally excludes objects already
missed by the baseline.

\thesisfigure{k_distance_pointrcnn}{PointRCNN loss by object range}{
Historical Car loss rates conditional on a qualifying baseline match
within each planar-distance bin. Lines connect the three archived
strata; they do not interpolate unmeasured performance.
PU-Net* and PU-EF retain the historical implementation qualifications.}{fig:k-distance-pr}
\thesisfigure{k_distance_centerpoint}{CenterPoint loss by object range}{
CenterPoint Car loss rates from the same historical audit.
The denominator is the baseline-matched count within each bin.
These descriptive strata are separate from KITTI Easy, Moderate,
and Hard difficulty definitions.}{fig:k-distance-cp}

\FloatBarrier
\FloatBarrier
\subsection{Whole-Scene Multi-Method Comparisons}
\label{sec:kitti-scene-visual}

\thesisfigure{k_scene_A}{Line A multi-method scene comparison}{
Lab frame 000590, historical Line~A E1 inputs. Gray points are retained
observed rows and blue points are generated rows. The shared crop is
0--60~m forward and $-18$--18~m lateral, all heights, with no display
subsampling. PU-Net* is the old faulty adapter; PU-EF is the lab
compatibility path.}{fig:k-scene-A}

\thesisfigure{k_scene_B}{Line B multi-method scene comparison}{
The same scene under historical Line~B construction. Measured points
now come from the sparse baseline. All four stored method outputs are
included under the same crop and colour convention as
Figure~\ref{fig:k-scene-A}.}{fig:k-scene-B}

Figures~\ref{fig:k-scene-A} and~\ref{fig:k-scene-B} re-render the local
lab visualisation arrays for frame 000590, including PDANS, PU-GCN,
PU-EdgeFormer, and PU-Net in both lines. The display window is fixed to
0--60~m forward and $-18$--18~m lateral, with no height filter.
All points inside this window are drawn. Measured rows are identified
by exact float32-row membership in the corresponding observed input,
rather than by assigning a method a presumed preservation property.

The separate colours reveal a distinction hidden by an all-blue dense
cloud: these historical E1 inputs retain observations at the stored-cloud
stage, but much of the displayed density comes from generated rows.
In Line~B, the measured component is the sparse input, so recovery
to the original file size does not restore the original measurements.
The figure's point counts refer to the display window, not to detector
input cardinality or the number of points inside labelled objects.

\FloatBarrier
\subsection{Local Surfaces and Measurement Preservation}
\label{sec:kitti-local-visual}

\thesisfigure{k_local_A}{Line A local multi-method surfaces}{
Common object-aligned crop for frame 005625, Car audit index 10.
All four outputs retain 316 measured rows in this window; generated
rows are shown separately. The dashed GT box and view are identical
across methods. Historical pipeline labels match Figure~\ref{fig:k-scene-A}.}{fig:k-local-A}

\thesisfigure{k_local_B}{Line B local multi-method surfaces}{
The same labelled Car after sparsification. All points within the
declared crop are shown; counts distinguish the 85 measured rows from
generated rows. Increased crop density alone does not establish
restored surface accuracy or improved detection.}{fig:k-local-B}

The whole-scene view cannot resolve the distribution around one object.
Figures~\ref{fig:k-local-A} and~\ref{fig:k-local-B} therefore use the
same Car in frame 005625 (audit GT index 10) and the same object-aligned
coordinates for every method. The crop extends 0.6~m beyond the box in
length and width and 0.4~m in height. These margins retain context
around the measured surface; crop counts are not restricted to the
interior of the ground-truth box.

The original crop contained 316 observed rows, compared with 85 after
sparsification. In Line~A, PDANS added 1,184 generated rows within this
window and PU-GCN added 953. In Line~B, the respective counts were
284 and 259. The visible point concentration and the incomplete
measurement support can therefore be assessed separately.
The dashed box specifies the labelled object's extent; it is not a
surface reconstruction target. Points outside it can be legitimate
background because the crop deliberately includes a margin.

\FloatBarrier
\subsection{Matched Detection Failures and Positive Counterexamples}
\label{sec:kitti-matched-cases}

\thesisfigure{k_case_pr_lost}{PointRCNN Car loss case}{Historical lab frame 000590, Line~A, PDANS, audit object index 2. The oblique view and BEV show the same effective-input crop. Predictions are read from the retained manifest; only qualifying matches are drawn. The class threshold is 0.70 3D IoU.}{fig:k-case-pr-lost}

The four examples below were already identified in the lab event audit.
They are selected diagnostic cases, not a random estimate of event
frequencies. Each compares the same labelled object before and after
upsampling using detector-effective input. Dashed boxes are ground
truth; solid boxes are saved qualifying predictions. Unqualified
predictions are not drawn, so an FN label means no qualifying match,
not necessarily an absence of proposals.

The PointRCNN Line~A PDANS example lost a Car that initially had a well-aligned prediction. The common local view permits inspection of altered neighbourhood support without confusing this with the global increase in stored point count.
Figure~\ref{fig:k-case-pr-lost} records 3D IoU 0.777 before upsampling and no qualifying match afterwards. The displayed effective-input crop contains 846 baseline points and 785 upsampled points.

\FloatBarrier

A different PointRCNN example provides a positive counterexample in Line~B. The recovered Car crosses the matching threshold despite fewer effective input points in this local crop. Local cardinality alone therefore cannot explain the sign of a detection transition.
Figure~\ref{fig:k-case-pr-recovered} records no qualifying match before upsampling and 3D IoU 0.705 afterwards. The displayed effective-input crop contains 312 baseline points and 271 upsampled points.

\thesisfigure{k_case_pr_recovered}{PointRCNN Car recovery case}{Historical lab frame 005625, Line~B, PDANS, audit object index 7. The oblique view and BEV show the same effective-input crop. Predictions are read from the retained manifest; only qualifying matches are drawn. The class threshold is 0.70 3D IoU.}{fig:k-case-pr-recovered}
\FloatBarrier

The CenterPoint Line~B PU-GCN Car is a complementary negative case: more retained crop points coexist with loss of a previously qualifying prediction. This bounds the claim that preserving measurements or adding local density is sufficient for detection.
Figure~\ref{fig:k-case-cp-lost} records 3D IoU 0.871 before upsampling and no qualifying match afterwards. The displayed effective-input crop contains 100 baseline points and 394 upsampled points.

\thesisfigure{k_case_cp_lost}{CenterPoint Car loss case}{Historical lab frame 005625, Line~B, PU-GCN, audit object index 10. The oblique view and BEV show the same effective-input crop. Predictions are read from the retained manifest; only qualifying matches are drawn. The class threshold is 0.70 3D IoU.}{fig:k-case-cp-lost}
\FloatBarrier

The CenterPoint Cyclist in the same Line~B PU-GCN frame provides a positive case under the class-specific 0.50 IoU threshold. It shows that recoveries can occur in the same representation and scene as Car losses; the full-validation event counts establish which direction dominates.
Figure~\ref{fig:k-case-cp-recovered} records no qualifying match before upsampling and 3D IoU 0.655 afterwards. The displayed effective-input crop contains 19 baseline points and 110 upsampled points.

\thesisfigure{k_case_cp_recovered}{CenterPoint Cyclist recovery case}{Historical lab frame 005625, Line~B, PU-GCN, audit object index 4. The oblique view and BEV show the same effective-input crop. Predictions are read from the retained manifest; only qualifying matches are drawn. The class threshold is 0.50 3D IoU.}{fig:k-case-cp-recovered}
\FloatBarrier

These cases discriminate several explanations that an AP table cannot.
A global fourfold expansion need not increase the number of points
entering a particular PointRCNN crop, and a larger CenterPoint crop
can still lose its matched box. Conversely, a positive transition does
not require more local effective points. The examples support analysing
where generated points enter the representation and how predictions
change, while leaving the causal contribution of geometry, sampling,
and learned features unresolved. The full-validation final matrix
remains the evidence for detector adaptation.

